\documentclass[12pt]{article}

% Configuración de página
\usepackage[margin=.85in]{geometry}

% Referencias
\usepackage{natbib}
\usepackage{nicefrac}

% Colores principales
\usepackage[table]{xcolor}
\definecolor{mainblue}{HTML}{2c3e50}
\definecolor{darkblue}{HTML}{003366}
\definecolor{lightgray}{HTML}{f7f9fb}

% Fuentes y tipografía
\usepackage[utf8]{inputenc}
\usepackage[T1]{fontenc}
\usepackage{lmodern}

% Configuración de títulos
\usepackage{titlesec}
\titleformat{\section}{\Large\bfseries}{\thesection}{1em}{}
\titleformat{\subsection}{\large\bfseries}{\thesubsection}{1em}{}
\titleformat{\subsubsection}{\normalsize\bfseries}{\thesubsubsection}{1em}{}

% Configuración del título principal
\usepackage{titling}
\pretitle{\begin{center}\LARGE\bfseries\vskip -3em}
\posttitle{\par\end{center}\vskip -3em}
\preauthor{\begin{center}\vskip -2.8em}
\postauthor{\par\end{center}\vskip -2.5em}
\predate{\begin{center}\vskip -2.6em}
\postdate{\par\end{center}\vskip -1em}

% Headers y footers
\usepackage{fancyhdr}
\pagestyle{fancy}
\fancyhf{}
\fancyhead[L]{Macroeconomía II: Quiz 1}
\fancyhead[R]{\thepage}
\renewcommand{\headrulewidth}{0.5pt}
\renewcommand{\headrule}{%
  \hbox to\headwidth{\leaders\hrule height \headrulewidth\hfill}
}

% Configuración del TOC
\usepackage{tocloft}
\renewcommand{\cftsecfont}{\bfseries}
\renewcommand{\cftsubsecfont}{}
\renewcommand{\cftsecpagefont}{}
\renewcommand{\cftsubsecpagefont}{}

% Espaciado y márgenes
\usepackage{setspace}
\onehalfspacing
\setlength{\parskip}{5pt}
\setlength{\parindent}{0pt}

% Configuración de enlaces
\usepackage{hyperref}
\hypersetup{
  colorlinks=true,
  linkcolor=black,
  urlcolor=darkblue,
  citecolor=darkblue
}

% Configuración de código para saltos de línea y tamaño
\usepackage{framed}
\usepackage{fancyvrb}
\usepackage{fvextra}
\usepackage{etoolbox}
\definecolor{codebg}{RGB}{250,250,255}
\definecolor{shadecolor}{named}{codebg}

\makeatletter
\@ifundefined{Shaded}{
  \newenvironment{Shaded}
    {\begin{snugshade}\scriptsize}
    {\end{snugshade}}
}{
  \renewenvironment{Shaded}
    {\begin{snugshade}\scriptsize}
    {\end{snugshade}}
}
\makeatother

\fvset{
  breaklines=true,
  breaksymbol=\textcolor{gray}{\tiny$\hookrightarrow$}
}

% Configuración de tcolorbox
\usepackage[most]{tcolorbox}
\tcbset{
  colframe=blue!30!black,
  colback=white,
  breakable,
  enhanced,
  boxrule=0.5mm,
  sharp corners,
  fonttitle=\bfseries,
  coltitle=white,
  colbacktitle=blue!30!black
}
\newtcolorbox{solution}{title=Solution}

\usepackage{tikz}
\usepackage{tikz-cd}
\usepackage{booktabs}
\usepackage{longtable}
\usepackage{array}
\usepackage{multirow}

% LaTeX
\usepackage{bbm}
\usepackage{pgfplots}
\pgfplotsset{compat=1.18}

% Math
\usepackage{amsthm,amsmath,amssymb}

\numberwithin{equation}{section}

\theoremstyle{definition}
\newtheorem{definition}{Definition}

\theoremstyle{plain}
\newtheorem{theorem}{Theorem}
\newtheorem{lemma}{Lemma}[section]
\newtheorem{proposition}{Proposition}[section]
\newtheorem{corollary}{Corollary}[theorem]

\theoremstyle{remark}
\newtheorem*{remark}{Remark}

\tcbuselibrary{theorems}

\tcolorboxenvironment{example}{
  enhanced,
  borderline north={0.5pt}{0pt}{black},
  borderline south={0.5pt}{0pt}{black},
  boxrule=0pt,
  frame hidden,
  before upper={\parindent15pt},
  footnote style
}

\tcolorboxenvironment{exercise}{
  enhanced,
  borderline north={0.5pt}{0pt}{black},
  borderline south={0.5pt}{0pt}{black},
  boxrule=0pt,
  frame hidden,
  before upper={\parindent15pt},
  footnote style
}

\tcolorboxenvironment{remark}{
  enhanced,
  borderline north={0.5pt}{0pt}{blue},
  borderline south={0.5pt}{0pt}{blue},
  boxrule=0pt,
  frame hidden,
  before upper={\parindent15pt},
  footnote style
}

\theoremstyle{definition}
\newtheorem{example}{Example}[section]
\newtheorem*{exercise}{Ex}

% Información del documento
\title{\large Macroeconomía II \vspace{0.5cm}}
\author{
  Centro de Investigación y Docencia Económicas\\
  \vspace{-0.05cm}
  Licenciatura en Economía -- Otoño 2026\\
  \vspace{-0.05cm}
  Quiz 1
}
\date{}

\begin{document}

\maketitle

\vspace{-.5cm}

\textbf{Matrícula:} \rule{15cm}{0.4pt}

\small
Anote su matrícula con pluma, \textbf{no incluya su nombre}.
Responda cada inciso de \textbf{forma detallada} y con
\textbf{letra legible}.

\vspace{.4cm}

\normalsize

\textbf{Ejercicio 1. (60 puntos)} Considere un hogar con preferencias estrictamente crecientes y estrictamente convexas sobre $(c,c')$, que recibe ingresos $(y,y')$ y paga impuestos de suma fija $(\tau,\tau')$. Puede ahorrar o endeudarse libremente a la tasa de interés real $r$. Suponga que ambos consumos son bienes normales.
Para cada uno de los siguientes casos, utilice un diagrama en el plano $(c,c')$ para analizar el efecto sobre el consumo presente, el consumo futuro y el ahorro actual.

\begin{enumerate}
  \item[(a)] El gobierno anuncia una \textbf{reducción} de los impuestos en $\Delta\tau<0$ únicamente en el periodo actual.

  \item[(b)] El gobierno anuncia una \textbf{reducción} de impuestos permanente: tanto $\tau$ como $\tau'$ se reducen en $\Delta\tau<0$.

  \item[(c)] Compare los resultados de los dos incisos. Explique por qué la
  respuesta del consumo depende de si el cambio en el ingreso disponible es
  transitorio o permanente.
\end{enumerate}

\textbf{Ejercicio 2. (40 puntos)} Suponga que un hogar resuelve el siguiente problema de consumo y ahorro en dos periodos:
\[
\max_{c_1,s,c_2}\;\frac{c_1^{1-\sigma}-1}{1-\sigma} + \beta \,\frac{c_2^{1-\sigma}-1}{1-\sigma} \qquad
\sigma>0
\]
sujeto a
\begin{align*}
s   &= s_0+y_1-\tau_1-c_1,\\
c_2 &= y_2-\tau_2+(1+r)s,
\end{align*}

donde $c_1$ es el consumo en el periodo 1; $c_2$ es el consumo en el periodo 2; $y_1$ es el ingreso del hogar en el periodo 1; $y_2$ es su ingreso en el periodo 2; $\tau_1$ y $\tau_2$ son los impuestos pagados en los periodos 1 y 2, respectivamente; $s_0$ es la riqueza inicial del hogar; $s$ es el ahorro elegido en el primer periodo; $r$ es la tasa de interés real; y $\beta$ es el factor subjetivo de descuento.

\begin{enumerate}
  \item[(a)] Obtenga la restricción presupuestaria intertemporal.
  \item[(b)] Obtenga la ecuación de Euler e interprétela económicamente.
\end{enumerate}

\end{document}
